% Consolidated 12-entropy-cycles from df0f579b9265792e2eecacd1f44bac71db6664c5; original TeX fragments preserved.

% === BEGIN 12-entropy-cycles-113-121.tex blob 93a453c7d59aa68db278424a4d179731cfea3cef ===
% Irodov Part II section 2.4, original 113--121, printed p.76, PDF p.85.
% Carnot, Otto, Brayton and mixed reversible ideal-gas engine cycles.
\thermoentropyproblem{113}{76}
在温度分别为 $T_h>T_c>0$ 的两个热源间工作的
可逆卡诺热机，哪种方式对提高热效率更有效：
将高温热源的温度升高相同的
$\Delta T>0$，
还是将低温热源的温度降低 $\Delta T$
（假设 $T_c>\Delta T$）？
\exsolution
卡诺效率
$\eta=1-T_c/T_h$。
提高高温热源温度得到
\[
\eta_h=1-\frac{T_c}{T_h+\Delta T},
\]
降低低温热源温度得到
\[
\eta_c=1-\frac{T_c-\Delta T}{T_h}.
\]
直接相减：
\[
\eta_c-\eta_h
=\frac{\Delta T}{T_h}
-\frac{T_c\Delta T}{T_h(T_h+\Delta T)}
=\frac{\Delta T(T_h+\Delta T-T_c)}
{T_h(T_h+\Delta T)}>0.
\]
故
\[
\boxed{\text{降低低温热源温度的同量 }\Delta T
\text{ 可得到更大的效率提高}}.
\]
\excheck 比较必须在仍有
$T_c-\Delta T>0$ 的物理范围中进行。

\thermoentropyproblem{114}{76}
氢气作为卡诺热机的工质。
若绝热膨胀阶段：
(a) 气体体积增加 $n=2.0$ 倍；
(b) 气体压强下降 $n=2.0$ 倍，
求各情况下的卡诺效率。
\exsolution
卡诺循环两等温温度为 $T_h,T_c$，
效率
$\eta=1-T_c/T_h$。
假设氢气为刚性双原子理想气体，
$\gamma=7/5$。
(a) 可逆绝热过程满足
$TV^{\gamma-1}=\text{常数}$，
若 $V_c/V_h=n$，则
$T_c/T_h=n^{-(\gamma-1)}$。
因此
\[
\boxed{\eta_a=1-n^{1-\gamma}
=1-2^{-2/5}\simeq0.24\approx25\%}.
\]
(b) 由可逆绝热关系
$T\propto p^{(\gamma-1)/\gamma}$，
$T_c/T_h=n^{-(\gamma-1)/\gamma}$。
故
\[
\boxed{\eta_b=1-n^{-(\gamma-1)/\gamma}
=1-2^{-2/7}\simeq0.18}.
\]
\excheck 两问给定的是不同的绝热压缩比，
不能把体积倍率与压力倍率直接混用。

\thermoentropyproblem{115}{76}
一台卡诺热机原先工作在两个热源之间，
其热机效率为 $\eta=10\%$。
若将同一装置倒转用作卡诺制冷机，
且冷热热源仍不变，
求制冷系数 $\varepsilon$。
\exsolution
对可逆卡诺循环，
热机效率
$\eta=1-T_c/T_h$。
反向制冷机吸收冷源热量 $Q_c$、
消耗外功 $W=Q_h-Q_c$，
制冷系数定义为
$\varepsilon=Q_c/W$。
卡诺关系
$Q_c/Q_h=T_c/T_h$，
因此
\[
\varepsilon=\frac{T_c}{T_h-T_c}
=\frac{T_c/T_h}{1-T_c/T_h}
=\frac{1-\eta}{\eta}.
\]
代入 $\eta=0.10$：
\[
\boxed{\varepsilon=\frac{0.90}{0.10}=9.0}.
\]
\excheck 制冷系数不是热机效率，
可以大于1，因为它衡量的是
搬运的冷源热量与外功之比。

\thermoentropyproblem{116}{76}
一理想气体经过由等温线和绝热线交替组成的
可逆热机循环（原图2.2）。
三段等温过程分别处于
$T_1>T_2>T_3$。
其中两个等温膨胀阶段的体积均以相同倍率
$r>1$ 增大。
求整个循环效率。
\exsolution
在温度 $T_1,T_2$ 上，
各发生一次等温膨胀，
理想气体等温吸热分别为
$Q_1=\nu RT_1\ln r$、
$Q_2=\nu RT_2\ln r$。
绝热连接段不交换热。
由于闭合循环内总熵变化为零，
低温 $T_3$ 的等温压缩须排出
$Q_3$ 满足
\[
\frac{Q_1}{T_1}+\frac{Q_2}{T_2}
-\frac{Q_3}{T_3}=0.
\]
故
\[
Q_3=2\nu RT_3\ln r.
\]
热机做净功
$W=Q_1+Q_2-Q_3$，
热效率为
\[
\boxed{\eta=\frac{W}{Q_1+Q_2}
=1-\frac{2T_3}{T_1+T_2}}.
\]
\excheck 倍率 $r$ 通过等温吸热表达式相消；
若忽略第二条中温等温线，
会错误地套用普通两热源卡诺公式。

\thermoentropyproblem{117}{76}
一理想热机循环由两条等容线和
两条可逆绝热线构成（Otto 循环）。
若体积比 $r=V_{\max}/V_{\min}=10$，
且工质为氮气，求效率。
\exsolution
在绝热压缩与绝热膨胀上，
$TV^{\gamma-1}=\text{常数}$。
令四状态的两等容吸放热
$Q_{\rm in}=\nu C_V(T_3-T_2)$、
$Q_{\rm out}=\nu C_V(T_4-T_1)$。
配对使用绝热关系
$T_2/T_1=T_3/T_4=r^{\gamma-1}$，
可得
\[
\frac{Q_{\rm out}}{Q_{\rm in}}
=r^{1-\gamma}.
\]
故循环效率
\[
\boxed{\eta=1-r^{1-\gamma}}.
\]
氮气取 $\gamma=7/5$，代入 $r=10$：
\[
\boxed{\eta=1-10^{-2/5}\simeq0.60=60\%}.
\]
\excheck 与两条\emph{等压}的循环不同，
这里的压缩倍率定义为体积倍率。

\thermoentropyproblem{118}{76}
一理想气体的可逆热机循环
由两条等压线和两条绝热线组成
（Brayton 循环）。
若循环中最高与最低压强之比为 $n$，
气体绝热指数为 $\gamma$，
求热效率。
\exsolution
设高低压强分别为
$p_h,p_l$，$p_h/p_l=n>1$。
在两绝热支路上
$T\propto p^{(\gamma-1)/\gamma}$，
故成对状态的温度比
\[
\frac{T_2}{T_1}=\frac{T_3}{T_4}
=n^{(\gamma-1)/\gamma}.
\]
高压等压支路吸热
$Q_{\rm in}=\nu C_p(T_3-T_2)$，
低压等压支路排热
$Q_{\rm out}=\nu C_p(T_4-T_1)$。
因此
\[
\frac{Q_{\rm out}}{Q_{\rm in}}
=n^{-(\gamma-1)/\gamma}.
\]
效率为
\[
\boxed{\eta=1-n^{-(\gamma-1)/\gamma}}.
\]
\excheck 必须保证循环方向为正功热机，
并使用正的吸热、排热量计算效率。

\thermoentropyproblem{119}{76}
一理想气体的顺时针循环由两条等容线
与两条等压线组成，
气体绝热指数为 $\gamma$。
若等容加热和等压膨胀时
气体的绝对温度都各自增加 $n$ 倍，
求循环效率。
\exsolution
取初态
$(p_1,V_1,T_1)$，
等容加热使压强变为 $np_1$、
温度变为 $nT_1$；
随后在高压 $np_1$ 下等压膨胀，
温度再增至 $n^2T_1$、
体积变为 $nV_1$。
剩下两支路为等容冷却和低压等压压缩。
循环矩形面积是净功：
\[
W=(np_1-p_1)(nV_1-V_1)
=p_1V_1(n-1)^2
=\nu RT_1(n-1)^2.
\]
吸热在等容加热和高压等压膨胀段：
\[
Q_{\rm in}=\nu C_VT_1(n-1)
+\nu C_pT_1n(n-1).
\]
因此
\[
\eta=\frac{R(n-1)}{C_V+nC_p}.
\]
用 $C_p=\gamma C_V$、
$R=(\gamma-1)C_V$ 化简：
\[
\boxed{\eta=\frac{(\gamma-1)(n-1)}
{1+\gamma n}}.
\]
\excheck 若 $n\to1$ 循环面积与功趋零，
效率也随之趋零。

\thermoentropyproblem{120}{76}
一理想气体经历由
(a) 等容、绝热、等温三支路；
(b) 等压、绝热、等温三支路组成的可逆热机循环。
假设等温支路始终处于\emph{最低温度}
$T_c$，循环的最高温度与最低温度之比为
$n=T_h/T_c>1$。
求两循环的效率。
\exsolution
(a) 等容加热将温度由 $T_c$ 提升至 $T_h$，
随后可逆绝热膨胀回到 $T_c$，
最后在低温等温压缩回初态。
等容吸热
$Q_{\rm in}=\nu C_V(T_h-T_c)
=\nu C_VT_c(n-1)$。
绝热支路给
$V_{\rm after}/V_{\rm before}
=n^{1/(\gamma-1)}$，
因此低温等温排热
$Q_{\rm out}=\nu RT_c\ln n/(\gamma-1)
=\nu C_VT_c\ln n$。
故
$\eta_a=1-\ln n/(n-1)$。
(b) 等压加热的吸热
$Q_{\rm in}=\nu C_pT_c(n-1)$。
绝热连接使低温等温压缩的体积比为
$n^{\gamma/(\gamma-1)}$，
从而
$Q_{\rm out}=\nu RT_c\gamma\ln n/(\gamma-1)
=\nu C_pT_c\ln n$。
于是两情况均有
\[
\boxed{\eta_a=\eta_b=1-\frac{\ln n}{n-1}}.
\]
\excheck 正循环必须沿吸热支路升温至
$T_h$；若把等温支路误当作最高温，
会得到第121题的不同结果。

\thermoentropyproblem{121}{76}
条件与上一题相同，
但等温支路改在整个循环的
\emph{最高温度} $T_h$ 上进行。
已知最高、最低温比为
$n=T_h/T_c>1$，
求两种循环各自的效率。
\exsolution
(a) 高温等温膨胀吸热，
随后等容冷却至 $T_c$，
再通过绝热压缩回高温起点。
绝热支路给
$V_{\rm high}/V_{\rm low}
=n^{1/(\gamma-1)}$，
高温等温吸热
$Q_{\rm in}=\nu RT_h\ln n/(\gamma-1)
=\nu C_VT_h\ln n$。
等容冷却排热
$Q_{\rm out}=\nu C_V(T_h-T_c)
=\nu C_VT_h(1-1/n)$，
故
$\eta_a=1-(n-1)/(n\ln n)$。
(b) 高温等温膨胀的体积比是
$n^{\gamma/(\gamma-1)}$，
所以
$Q_{\rm in}=\nu C_pT_h\ln n$；
等压冷却排热
$Q_{\rm out}=\nu C_pT_h(1-1/n)$。
两种情况同为
\[
\boxed{\eta_a=\eta_b=1-\frac{n-1}{n\ln n}}.
\]
\excheck $n>1$ 时
$\ln n>(n-1)/n$，
故效率为正且小于1。
% === END 12-entropy-cycles-113-121.tex ===

% === BEGIN 12-entropy-cycles-122-129.tex blob c2f46d76dd9f345bcdd9928b784c048117abeba6 ===
% Irodov, section 2.4 problems 122--129, printed pp.76--77,
% PDF pp.85--86. Fig 2.3 original inspected for #127.
\thermoentropyproblem{122}{76}
理想气体经历由\emph{一条等温线、
一条多方线及一条绝热线}组成的可逆热机循环，
等温阶段位于循环的最高温度 $T_h$。
求热效率，并以循环最高温度与最低温度之比
$\tau=T_h/T_c>1$ 表示。
\exsolution
设多方指数为 $n$，与绝热指数 $\gamma$ 不同，
并取其支路在冷却时热容量为正，
使循环沿正功方向运行。
沿两条非等温支路，
可逆绝热满足
$TV^{\gamma-1}=\text{常数}$，
多方满足
$TV^{n-1}=\text{常数}$。
设高温等温膨胀的体积倍率为 $r$。
由两个连接段分别跨越相同的温比 $\tau$，
\[
\ln r=
\ln\tau\left(\frac1{\gamma-1}-\frac1{n-1}\right).
\]
在高温等温支路吸热
\[
Q_{\rm in}=\nu RT_h\ln r
=\nu RT_h\ln\tau
\left(\frac1{\gamma-1}-\frac1{n-1}\right).
\]
多方冷却支路的摩尔热容为
\[
C_n=C_V+\frac R{1-n}
=R\left(\frac1{\gamma-1}-\frac1{n-1}\right).
\]
它的排热量为
$Q_{\rm out}=\nu C_n(T_h-T_c)$，
绝热支路无热交换。
因此
\[
\frac{Q_{\rm out}}{Q_{\rm in}}
=\frac{T_h-T_c}{T_h\ln\tau}
=\frac{\tau-1}{\tau\ln\tau},
\]
\[
\boxed{\eta_{\rm cycle}
=1-\frac{\tau-1}{\tau\ln\tau}}.
\]
\excheck 本版第122题题面未给出温度倍率的数值，
所以没有可计算的数值效率。
这里显式定义 $\tau=T_h/T_c$，
并假设热源过程方向与可逆热机一致；
与第121题的函数形式相同，
不是给题目虚构了一个数值倍率。

\thermoentropyproblem{123}{77}
一理想气体的顺时针正功循环
由一条可逆绝热、一条等压和一条等容线组成。
气体绝热指数为 $\gamma$。
如果绝热阶段体积
(a) 增大 $n$ 倍；
(b) 缩小 $n$ 倍，
求对应的热效率。
\exsolution
(a) 令绝热膨胀初态为
$A:(p_1,V_1,T_1)$，
末态
$B:(p_1n^{-\gamma},nV_1,T_1n^{1-\gamma})$。
然后在低压下等压压缩至
$C:(p_1n^{-\gamma},V_1,T_1n^{-\gamma})$，
最后等容加热回 $A$。
吸热发生于 $C\to A$，
排热发生于 $B\to C$：
\[
Q_{\rm in}=\nu C_VT_1(1-n^{-\gamma}),
\quad
Q_{\rm out}=\nu C_pT_1n^{-\gamma}(n-1).
\]
所以
\[
\boxed{\eta_a=
1-\frac{\gamma(n-1)}{n^\gamma-1}}.
\]
(b) 令绝热压缩初态为
$A:(p_1,nV_1,T_1)$，
压缩至
$B:(p_1n^\gamma,V_1,T_1n^{\gamma-1})$。
然后在高压下等压膨胀至
$C:(p_1n^\gamma,nV_1,T_1n^\gamma)$，
最后等容冷却回 $A$。
吸热在 $B\to C$，
排热在 $C\to A$：
\[
Q_{\rm in}=\nu C_pT_1n^{\gamma-1}(n-1),
\quad Q_{\rm out}=\nu C_VT_1(n^\gamma-1).
\]
故
\[
\boxed{\eta_b=
1-\frac{n^\gamma-1}
{\gamma n^{\gamma-1}(n-1)}}.
\]
\excheck 必须明确绝热阶段是膨胀还是压缩，
否则会混淆两条正功循环的吸放热支路。

\thermoentropyproblem{124}{77}
由等温、等压及等容三段组成的理想气体
顺时针热机循环，绝热指数为 $\gamma$。
若等温支路上的体积
(a) 增大 $n$ 倍；(b) 缩小 $n$ 倍，
求循环效率。
\exsolution
(a) 从状态
$A:(p_1,V_1,T_1)$ 等温膨胀到
$B:(p_1/n,nV_1,T_1)$。
然后在低压 $p_1/n$ 下等压压缩至
$C:(p_1/n,V_1,T_1/n)$，
最后等容加热回 $A$。
吸热来自高温等温与末段等容加热：
\[
Q_{\rm in}=\nu RT_1\ln n+
\nu C_VT_1(1-1/n),
\]
排热来自等压压缩：
$Q_{\rm out}=\nu C_pT_1(1-1/n)$。
计算效率：
\[
\boxed{\eta_a=
1-\frac{\gamma(n-1)}
{n-1+(\gamma-1)n\ln n}}.
\]
(b) 从
$A:(p_1/n,nV_1,T_1)$
等温压缩到
$B:(p_1,V_1,T_1)$。
然后在高压 $p_1$ 下等压膨胀至
$C:(p_1,nV_1,nT_1)$，
最后等容冷却回 $A$。
吸热仅在等压膨胀：
$Q_{\rm in}=\nu C_pT_1(n-1)$；
排热为等温压缩的
$\nu RT_1\ln n$ 和等容冷却的
$\nu C_VT_1(n-1)$。
所以
\[
\boxed{\eta_b=
1-\frac{n-1+(\gamma-1)\ln n}
{\gamma(n-1)}}.
\]
\excheck (a) 的热输入有\emph{两条支路}，
而 (b) 的热输出有两条支路，
应逐项统计而非只套普通卡诺效率。

\thermoentropyproblem{125}{77}
由两条等容线和两条等温线构成理想气体
正功循环，工质绝热指数为 $\gamma$。
整个循环最大体积与最小体积之比为
$\upsilon>1$，最高温与最低温之比为
$\tau>1$。求\emph{未配置回热器}的循环效率。
\exsolution
取高温等温膨胀段温度 $T_h$，
低温等温压缩段温度 $T_c=T_h/\tau$，
体积倍率
$\upsilon=V_{\max}/V_{\min}$。
循环净功等于两个等温功之差：
\[
W=\nu R(T_h-T_c)\ln\upsilon.
\]
吸热不仅来自高温等温膨胀，
还来自在最小体积处的等容加热：
\[
Q_{\rm in}=\nu RT_h\ln\upsilon+
\nu C_V(T_h-T_c).
\]
所以
\[
\eta=
\frac{(T_h-T_c)\ln\upsilon}
{T_h\ln\upsilon+(C_V/R)(T_h-T_c)}.
\]
用 $C_V/R=1/(\gamma-1)$ 及
$T_h/T_c=\tau$：
\[
\boxed{\eta=
\frac{(\tau-1)\ln\upsilon}
{\tau\ln\upsilon+(\tau-1)/(\gamma-1)}}.
\]
\excheck 若理想回热器完全回收等容冷却的热，
外部吸热计算不同；原书本题未设回热器，
所以必须保留等容加热项。

\thermoentropyproblem{126}{77}
理想气体经由两条等压线和两条等温线组成
顺时针热机循环，绝热指数为 $\gamma$。
循环最高和最低压强之比为 $n>1$，
最高与最低温之比为 $\tau>1$。
求\emph{无回热}循环的效率。
\exsolution
高温等温膨胀的压强从高到低，
低温等温压缩的压强从低到高，
等温压强比均为 $n$，
故循环净功
\[
W=\nu R(T_h-T_c)\ln n.
\]
外界向气体供热在
高温等温膨胀和高压等压加热两段：
\[
Q_{\rm in}=\nu RT_h\ln n+
\nu C_p(T_h-T_c).
\]
于是
\[
\eta=\frac{(T_h-T_c)\ln n}
{T_h\ln n+(C_p/R)(T_h-T_c)}.
\]
使用
$C_p/R=\gamma/(\gamma-1)$，
\[
\boxed{\eta=
\frac{(\tau-1)\ln n}
{\tau\ln n+\gamma(\tau-1)/(\gamma-1)}}.
\]
\excheck 热量计入外界供热，
不假设存在理想回热器；
若忽略高压等压加热，效率将被高估。

\thermoentropyproblem{127}{77}
原书图2.3的顺时针 $p$–$V$ 三角形循环中，
一边为过原点的直线 $p=\alpha V$，
其余两边分别为等容线与等压线。
若循环最高温和最低温之比为
$\tau>1$，工质的绝热指数为 $\gamma$，
求效率。
\exsolution
取最低温点
$A=(V_1,p_1)$，最高温点
$B=(V_2,p_2)$，
二者在直线 $p=\alpha V$ 上，
因此若 $\lambda=V_2/V_1>1$，
则 $p_2/p_1=\lambda$，
$T_B/T_A=\lambda^2$，
故 $\lambda=\sqrt\tau$。
按原图方向，
$A\to B$ 沿 $p\propto V$ 吸热，
$B\to C$ 在 $V_2$ 等容冷却，
$C\to A$ 在 $p_1$ 等压压缩。
沿 $p\propto V$ 的摩尔热容为
$C_{AB}=R(\gamma+1)/[2(\gamma-1)]$，
于是
\[
Q_{\rm in}
=\nu\frac{R(\gamma+1)}{2(\gamma-1)}
(T_B-T_A)
=\frac{\gamma+1}{2(\gamma-1)}
p_1V_1(\lambda^2-1).
\]
闭合循环净功为三角形面积：
\[
W=\frac12(p_2-p_1)(V_2-V_1)
=\frac12p_1V_1(\lambda-1)^2.
\]
相除：
\[
\boxed{\eta=
\frac{\gamma-1}{\gamma+1}
\frac{\sqrt\tau-1}{\sqrt\tau+1}}.
\]
\excheck 原图箭头是\emph{顺时针}；
若颠倒运动方向则为需输入外功的逆循环，
不能把负功解释成热机效率。

\thermoentropyproblem{128}{77}
利用 Clausius 不等式证明：
在具有相同最高温度 $T_{\max}$ 和
最低温度 $T_{\min}$ 的热机循环中，
没有任何循环能比在这两个温度间工作的
可逆卡诺循环效率更高。
\exsolution
令某循环吸热段从环境吸收总热量 $Q_{\rm in}>0$，
放热段向环境放出热量模 $Q_{\rm out}>0$。
Clausius 不等式
\[
\oint\frac{\delta Q}{T}\le0
\]
可拆为吸热与放热两部分。
因为所有吸热发生的温度
$T\le T_{\max}$，
所以
$\int_{\rm in}\delta Q/T\ge Q_{\rm in}/T_{\max}$。
所有放热发生的温度
$T\ge T_{\min}$，
所以
$\int_{\rm out}\delta Q_{\rm out}/T
\le Q_{\rm out}/T_{\min}$。
因此
\[
0\ge
\oint\frac{\delta Q}{T}
=\int_{\rm in}\frac{\delta Q}{T}
-\int_{\rm out}\frac{\delta Q_{\rm out}}T
\ge\frac{Q_{\rm in}}{T_{\max}}
-\frac{Q_{\rm out}}{T_{\min}}.
\]
于是
$Q_{\rm out}/Q_{\rm in}\ge T_{\min}/T_{\max}$，
所以
\[
\boxed{\eta=\frac{W}{Q_{\rm in}}
=1-\frac{Q_{\rm out}}{Q_{\rm in}}
\le1-\frac{T_{\min}}{T_{\max}}
=\eta_{\rm Carnot}}.
\]
\excheck 等号要求理想可逆且吸热全在
$T_{\max}$、放热全在 $T_{\min}$
进行；一般不可逆循环严格低于此上界。

\thermoentropyproblem{129}{77}
利用卡诺定理证明：对由温度 $T$ 和体积 $V$
决定状态的均匀简单可压缩物质，
其内能 $U(T,V)$ 与状态方程 $p(T,V)$ 满足
\[
\left(\frac{\partial U}{\partial V}\right)_T
=T\left(\frac{\partial p}{\partial T}\right)_V-p.
\]
提示：考虑 $p$–$V$ 图上的无穷小卡诺循环。
\exsolution
在温度 $T$ 和 $T+\mathrm dT$ 之间
做无穷小可逆卡诺循环，
高温等温阶段的体积变化为 $\mathrm dV$。
卡诺定理给出效率的一阶式
\[
\frac{\delta W}{\delta Q_h}=\frac{\mathrm dT}{T}.
\]
在无穷小循环极限，
等温线之间的压差为
$\mathrm dp=(\partial p/\partial T)_V\,\mathrm dT$，
故循环净功的一阶项为
\[
\delta W=\left(\frac{\partial p}{\partial T}\right)_V
\mathrm dT\,\mathrm dV.
\]
高温等温支路吸热由第一定律
\[
\delta Q_h=\mathrm dU+p\,\mathrm dV
=\left[\left(\frac{\partial U}{\partial V}\right)_T+p\right]
\mathrm dV.
\]
代入卡诺效率并约去微元：
\[
\frac{(\partial p/\partial T)_V\,\mathrm dT}
{(\partial U/\partial V)_T+p}
=\frac{\mathrm dT}{T}.
\]
因此
\[
\boxed{\left(\frac{\partial U}{\partial V}\right)_T
=T\left(\frac{\partial p}{\partial T}\right)_V-p}.
\]
\excheck 一阶近似忽略的项是比
$\mathrm dT\,\mathrm dV$ 更高阶的微小面积贡献；
所得结果也可由 Helmholtz 自由能的
Maxwell 关系独立复核。
% === END 12-entropy-cycles-122-129.tex ===

% === BEGIN 12-entropy-cycles-130-139.tex blob a700f054c43e0d0c77bc7bd393fd43d20ae6c071 ===
% Irodov Part II section 2.4, original 130--139, printed pp.77--78
% (PDF pp.86--87). All entropy expressions are state functions.
\thermoentropyproblem{130}{77}
一摩尔二氧化碳气体（按理想气体处理）
绝对温度升高至原来的
$n=2.0$ 倍。
求气体熵增：
(a) 等容加热；
(b) 等压加热。
\exsolution
理想气体摩尔熵微分为
\[
\mathrm dS=C_V\frac{\mathrm dT}{T}
+R\frac{\mathrm dV}{V}
=C_p\frac{\mathrm dT}{T}-R\frac{\mathrm dp}{p}.
\]
(a) 等容 $V_2=V_1$，
\[
\boxed{\Delta S_V=C_V\ln\frac{T_2}{T_1}
=\frac{R}{\gamma-1}\ln n}.
\]
(b) 等压 $p_2=p_1$，
\[
\boxed{\Delta S_p=C_p\ln\frac{T_2}{T_1}
=\frac{\gamma R}{\gamma-1}\ln n}.
\]
按本版答案使用二氧化碳
$\gamma\simeq1.30$ 的模型值，
$n=2$，
\[
\boxed{\Delta S_V\simeq19\,\mathrm{J/K}},\qquad
\boxed{\Delta S_p\simeq25\,\mathrm{J/K}}.
\]
\excheck 二氧化碳的热容温度依赖较明显，
上述数值采用 $\gamma$ 近似常数；
更高精度应对 $C(T)/T$ 积分。

\thermoentropyproblem{131}{77}
$\nu=4.0$ mol 理想气体
在等温膨胀中熵增加
$\Delta S=23\,\mathrm{J/K}$。
求气体体积增大多少倍 $n=V_2/V_1$。
\exsolution
一摩尔理想气体等温熵变为
$R\ln(V_2/V_1)$，
故 $\nu$ 摩尔时
\[
\Delta S=\nu R\ln n.
\]
取指数：
\[
\boxed{n=\exp\left(\frac{\Delta S}{\nu R}\right)}.
\]
代入
\[
\boxed{n=\exp\left(\frac{23}{4(8.314)}\right)
\simeq2.0}.
\]
\excheck 体积比无量纲，若 $\Delta S>0$
则 $n>1$，符合等温自由膨胀或可逆膨胀的终态。

\thermoentropyproblem{132}{77}
$\nu=2.0$ mol 理想气体先等容冷却，
再等压膨胀使气体温度回到初值。
若整个过程中气压总共降低
$n=3.3$ 倍，
求系统熵变。
\exsolution
设初态 $(p_1,V_1,T_1)$。
第一段等容冷却到
$p_2=p_1/n$，
温度降至
$T_2=T_1/n$。
第二段等压升温至 $T_3=T_1$，
所以最终体积
$V_3=nV_1$，
最终压强仍为 $p_1/n$。
熵是状态函数，
由初末温度相同得
\[
\Delta S=\nu C_V\ln(T_3/T_1)
+\nu R\ln(V_3/V_1)
=\nu R\ln n.
\]
代入：
\[
\boxed{\Delta S=(2.0)R\ln(3.3)
\simeq20\,\mathrm{J/K}}.
\]
\excheck 中间过程先放热后吸热，
但熵变由初末状态确定，
不必强行把实际传热视为一个等温过程。

\thermoentropyproblem{133}{78}
氦气质量 $m=1.7\,\mathrm g$，
先进行可逆绝热膨胀，
体积增大 $n=3.0$ 倍；
随后在新压强下等压压缩，
直到气体总体积恢复初值。
求整个过程的气体熵变。
\exsolution
氦气是单原子理想气体，
$\gamma=5/3$、
$C_V=3R/2$、
$C_p=5R/2$。
设初态
$(p_1,V_1,T_1)$，
绝热膨胀后
$V_2=nV_1$，
$p_2=p_1n^{-\gamma}$。
后续等压压缩至
$V_3=V_1$，
所以末态
$p_3=p_2=p_1n^{-\gamma}$，
$T_3=T_1n^{-\gamma}$。
初末体积相同，因此
\[
\Delta S=\nu C_V\ln(T_3/T_1)
=-\nu\gamma C_V\ln n
=-\nu C_p\ln n.
\]
摩尔数
$\nu=m/M_{\rm He}=1.7/4.0=0.425$ mol，
故
\[
\boxed{\Delta S=
-\frac m{M_{\rm He}}\frac{5R}{2}\ln3
\simeq-10\,\mathrm{J/K}}.
\]
\excheck 气体系统熵可以在\emph{向环境放热}时降低；
第二定律要求整体包括环境的熵变化非负，
不要求任一子系统熵永远增加。

\thermoentropyproblem{134}{78}
$\nu=2.0$ mol 理想气体，
绝热指数 $\gamma=1.30$。
某过程中体积增大到
$\alpha=2.0$ 倍，而压强减为
$1/\beta$，$\beta=3.0$。
求气体熵变。
\exsolution
由状态方程
$T\propto pV$，
温度倍率为
$T_2/T_1=(p_2V_2)/(p_1V_1)=\alpha/\beta$。
对理想气体
\[
\Delta S=\nu C_V\ln(T_2/T_1)
+\nu R\ln(V_2/V_1).
\]
代入 $C_V=R/(\gamma-1)$：
\[
\Delta S=\frac{\nu R}{\gamma-1}
\left[\ln\left(\frac\alpha\beta\right)
+(\gamma-1)\ln\alpha\right].
\]
故
\[
\boxed{\Delta S=
\frac{\nu R}{\gamma-1}
\left(\gamma\ln\alpha-\ln\beta\right)}.
\]
给定数值：
\[
\boxed{\Delta S=
\frac{2R}{0.30}
[1.30\ln2-\ln3]\simeq-11\,\mathrm{J/K}}.
\]
\excheck 即使体积增大，若降温和降压足够显著，
系统熵变仍可能为负；
不应仅凭膨胀断定熵增加。

\thermoentropyproblem{135}{78}
容器1、2各含
$\nu=1.2$ mol 理想氦气。
两容器的容积比
$V_2/V_1=\alpha=2.0$，
绝对温度比
$T_1/T_2=\beta=1.5$。
求两种状态下氦气熵的差
$S_2-S_1$。
\exsolution
理想气体在两个平衡状态间的熵差为
\[
S_2-S_1=
\nu C_V\ln\frac{T_2}{T_1}
+\nu R\ln\frac{V_2}{V_1}.
\]
氦气为单原子气体，
$C_V=3R/2$，于是
\[
\Delta S=\nu R
\left(-\frac32\ln\beta+\ln\alpha\right).
\]
代入
$\nu=1.2,\alpha=2,\beta=1.5$：
\[
\boxed{S_2-S_1=
1.2R\left(\ln2-\frac32\ln1.5\right)
\simeq0.85\,\mathrm{J/K}\approx1.0\,\mathrm{J/K}}.
\]
\excheck 两容器各自都含相同物质的量，
可将它们看作相同气体的两个状态比较；
无需考虑不同气体混合的混合熵。

\thermoentropyproblem{136}{78}
一摩尔绝热指数为 $\gamma$ 的理想气体
经多方过程 $pV^n=\text{常数}$，
末温为初温的 $\tau$ 倍。
求该过程的熵变。
\exsolution
多方过程
$TV^{n-1}=\text{常数}$，
故
\[
\ln(V_2/V_1)=-\frac{\ln\tau}{n-1}.
\]
理想气体熵变化为
\[
\Delta S=C_V\ln(T_2/T_1)
+R\ln(V_2/V_1).
\]
代入 $C_V=R/(\gamma-1)$ 得
\[
\Delta S=\frac R{\gamma-1}\ln\tau
-\frac R{n-1}\ln\tau.
\]
整理：
\[
\boxed{\Delta S=
\frac{R(n-\gamma)}{(n-1)(\gamma-1)}
\ln\tau}.
\]
\excheck $n=\gamma$ 是可逆绝热过程，
$\Delta S=0$；
$n=1$ 是等温多方过程，此时 $\tau=1$，
若两端都是1则须单独用体积比计算熵变，
不能简单除以 $n-1$。

\thermoentropyproblem{137}{78}
$\nu=2.0$ mol 氩气在一次准静态膨胀中
压强始终与体积成正比，
$p=\alpha V$。
若气体体积增大到原来的
$r=2.0$ 倍，
求气体熵变。
\exsolution
由 $pV=\nu RT$ 和
$p\propto V$，
$T\propto V^2$，
因此
$T_2/T_1=r^2$。
单原子氩气的定容摩尔热容
$C_V=3R/2$，于是
\[
\Delta S=\nu C_V\ln(r^2)+\nu R\ln r
=\nu(2C_V+R)\ln r.
\]
因 $2C_V+R=4R$，
\[
\boxed{\Delta S=4\nu R\ln r
=8R\ln2\simeq46\,\mathrm{J/K}}.
\]
\excheck 压强与体积正比是该题的\emph{特定路径}，
不能把这一膨胀误当作等压或等温过程。

\thermoentropyproblem{138}{78}
一摩尔绝热指数为 $\gamma$ 的理想气体
沿 $p=p_0-\alpha V$
准静态变化，其中 $p_0,\alpha>0$。
求气体熵随过程变化时达到极大值所对应的体积。
\exsolution
由
$RT=pV=p_0V-\alpha V^2$，
沿过程
\[
\frac{\mathrm dT}{\mathrm dV}
=\frac{p_0-2\alpha V}{R}.
\]
熵微分为
$\mathrm dS=C_V\,\mathrm dT/T+
R\,\mathrm dV/V$，
所以
\[
\frac{\mathrm dS}{\mathrm dV}
=\frac{C_V(p_0-2\alpha V)}
{p_0V-\alpha V^2}+\frac R V.
\]
令导数为零并通分，
\[
C_V(p_0-2\alpha V)+R(p_0-\alpha V)=0.
\]
即
\[
(C_V+R)p_0-(2C_V+R)\alpha V=0.
\]
由于 $C_p=\gamma C_V$、
$2C_V+R=(\gamma+1)C_V$，
\[
\boxed{V_{\rm max}=\frac{\gamma p_0}{\alpha(1+\gamma)}}.
\]
\excheck 物理定义域是
$0<V<p_0/\alpha$；
因为 $\gamma/(\gamma+1)<1$，
极值位于允许区间内，且为熵最大值。

\thermoentropyproblem{139}{78}
一摩尔理想气体沿某过程满足
\[
S(T)=aT+C_V\ln T,
\]
其中 $a>0$、
$C_V$ 为气体定容摩尔热容。
已知体积为 $V_0$ 时温度为 $T_0$，
求温度随体积 $T(V)$ 的变化规律。
\exsolution
由理想气体的平衡态熵表达式
\[
S(T,V)=C_V\ln T+R\ln V+\text{常数},
\]
将其与题给 $S(T)=aT+C_V\ln T$ 比较，
两边的 $C_V\ln T$ 抵消后，
有
\[
R\ln V=aT+\text{常数}.
\]
用初始条件 $V=V_0,T=T_0$ 消去常数：
\[
R\ln\frac V{V_0}=a(T-T_0).
\]
因此
\[
\boxed{T(V)=T_0+\frac Ra\ln\frac V{V_0}}.
\]
\excheck 对 $a>0$，体积增大时温度上升；
模型需限制在 $T(V)>0$ 的物理定义域内。
% === END 12-entropy-cycles-130-139.tex ===

% === BEGIN 12-entropy-cycles-140-145.tex blob bc879cceaa3713081bf9899fcbafd496014313f7 ===
% Irodov, Part II section 2.4, original problems 140--145,
% printed pp.78--79, PDF pp.87--88.
\thermoentropyproblem{140}{78}
一摩尔范德瓦耳斯气体沿等温过程，
由体积 $V_1$ 变化到 $V_2$。
已知范德瓦耳斯参数 $a,b$，
求气体熵的变化 $\Delta S$。
\exsolution
一摩尔范德瓦耳斯状态方程为
\[
p=\frac{RT}{V-b}-\frac a{V^2}.
\]
对任一单相平衡物质，
Maxwell 关系给
\[
\left(\frac{\partial S}{\partial V}\right)_T
=\left(\frac{\partial p}{\partial T}\right)_V.
\]
在范德瓦耳斯气体中
$(\partial p/\partial T)_V=R/(V-b)$。
因此沿等温路径：
\[
\Delta S=\int_{V_1}^{V_2}
\frac R{V-b}\,\mathrm dV
=\boxed{R\ln\frac{V_2-b}{V_1-b}}.
\]
\excheck 引力修正参数 $a$ 在
$\partial p/\partial T$ 中被消去，
等温熵变只显式依赖排除体积参数 $b$；
条件是 $V_1,V_2>b$。

\thermoentropyproblem{141}{78}
一摩尔范德瓦耳斯气体从
初态 $(V_1,T_1)$ 到
末态 $(V_2,T_2)$，
其定容摩尔热容 $C_V$ 已知并视为常数。
求熵变。
\exsolution
将熵视为状态函数 $S(T,V)$，
其全微分为
\[
\mathrm dS=
\left(\frac{\partial S}{\partial T}\right)_V\mathrm dT+
\left(\frac{\partial S}{\partial V}\right)_T\mathrm dV.
\]
由定容热容定义
$(\partial S/\partial T)_V=C_V/T$，
由 Maxwell 关系与范德瓦耳斯状态方程
$(\partial S/\partial V)_T=R/(V-b)$，
因此
\[
\mathrm dS=\frac{C_V}{T}\,\mathrm dT+
\frac R{V-b}\,\mathrm dV.
\]
沿任意连接两状态的可逆路径积分：
\[
\boxed{\Delta S=C_V\ln\frac{T_2}{T_1}
+R\ln\frac{V_2-b}{V_1-b}}.
\]
\excheck 由于熵是状态函数，
实际过程即使不可逆，也可用连接初末平衡态的
假想可逆路径计算熵差。

\thermoentropyproblem{142}{78}
在很低温的区间中，
某晶体的总热容满足
$C(T)=aT^3$，其中 $a$ 为常数。
求晶体在此温区的熵
$S(T)$，并说明第三定律边界。
\exsolution
对可逆温度变化，
熵微分
$\mathrm dS=\delta Q_{\rm rev}/T
=C(T)\mathrm dT/T$。
代入 $C(T)=aT^3$：
\[
\mathrm dS=aT^2\,\mathrm dT.
\]
从参考温度 $T_0$ 积分：
\[
S(T)-S(T_0)
=\frac a3(T^3-T_0^3).
\]
若晶体基态唯一且满足第三定律
$S(0)=0$，可令 $T_0=0$，
于是
\[
\boxed{S(T)=\frac a3T^3}.
\]
\excheck 若存在残余熵，
一般表达式为
$S(T)=S(0)+aT^3/3$，
不能未经条件说明就强制设 $S(0)=0$。

\thermoentropyproblem{143}{78}
质量 $m=3.0\,\mathrm{kg}$ 的铝棒
从 $T_1=300\,\mathrm K$ 加热到
$T_2=600\,\mathrm K$。
在此温区其质量比热容为
$c(T)=a+bT$，其中
$a=0.77\,\mathrm{J/(g\,K)}$、
$b=0.46\,\mathrm{mJ/(g\,K^2)}$。
求棒的熵增。
\exsolution
对温度相关热容，
应积分 $\mathrm dS=mc(T)\mathrm dT/T$，
不能先求平均热容再乘 $\ln(T_2/T_1)$。
所以
\[
\Delta S=m\int_{T_1}^{T_2}
\frac{a+bT}{T}\,\mathrm dT
=m\left[a\ln\frac{T_2}{T_1}
+b(T_2-T_1)\right].
\]
取 $m=3000\,\mathrm g$、
$b=0.00046\,\mathrm{J/(g\,K^2)}$，
代入：
\[
\boxed{\Delta S=3000
[0.77\ln2+0.00046(300)]
\,\mathrm{J/K}
\simeq2.0\times10^3\,\mathrm{J/K}}.
\]
\excheck $b$ 的题给单位是
$\mathrm{mJ/(g\,K^2)}$，
须先乘 $10^{-3}$ 化为焦耳单位。

\thermoentropyproblem{144}{79}
某物质在一条过程中，
温度与熵的关系为
$T=aS^n$，其中 $a,n$ 均为常数。
求该过程的热容 $C(S)$，
并说明热容为负的条件。
\exsolution
热容由可逆热量定义：
\[
C=\frac{\delta Q_{\rm rev}}{\mathrm dT}
=T\frac{\mathrm dS}{\mathrm dT}.
\]
由 $T=aS^n$，
在 $S>0$、$n\ne0$ 的定义域：
\[
\frac{\mathrm dT}{\mathrm dS}
=anS^{n-1}=\frac{nT}{S}.
\]
所以
\[
\boxed{C(S)=\frac{T}{\mathrm dT/\mathrm dS}
=\frac{S}{n}}.
\]
对于正的绝对熵 $S>0$，
\[
\boxed{C<0\iff n<0}.
\]
\excheck 若 $n=0$，温度不随熵变化，
通常的有限热容表达式不适用；
热容为负的情形也必须结合体系
热力学稳定性和所处约束来理解。

\thermoentropyproblem{145}{79}
某物质沿一条\emph{热容恒为 $C$}
的多方路径变化。
当 $T=T_0$ 时熵为 $S_0$。
求温度作为熵的函数 $T(S)$，
并定性描绘 $C>0$ 与 $C<0$ 的图线。
\exsolution
由热容与熵关系
$C=T\,\mathrm dS/\mathrm dT$，
分离变量：
\[
\mathrm dS=C\,\frac{\mathrm dT}{T}.
\]
从 $(T_0,S_0)$ 积分到 $(T,S)$，
\[
S-S_0=C\ln(T/T_0).
\]
因此
\[
\boxed{T(S)=T_0
\exp\left(\frac{S-S_0}{C}\right)}.
\]
当 $C>0$，$T$ 随熵 $S$
单调指数增加；当 $C<0$，
$T$ 随熵单调指数下降，且二者均通过
点 $(S_0,T_0)$。
\begin{center}
\begin{tikzpicture}[scale=0.7,>=Stealth]
\draw[->] (0,0)--(4.8,0) node[right] {$S$};
\draw[->] (0,0)--(0,3.5) node[above] {$T$};
\draw[thick,domain=0.3:4.0,samples=45]
plot (\x,{0.40*exp(0.48*\x)});
\node at (4.1,3.2) {$C>0$};
\draw[thick,dashed,domain=0.3:4.2,samples=45]
plot (\x,{2.8*exp(-0.42*\x)});
\node at (3.65,0.38) {$C<0$};
\end{tikzpicture}
\end{center}
\excheck 熵沿路径增加不必然表示温度上升，
这一点由 $C$ 的符号决定；
图示只表达单调性，不按特定材料数值比例绘制。
% === END 12-entropy-cycles-140-145.tex ===

% === BEGIN 12-entropy-cycles-146-159.tex blob 2dbc4ebb76f4a83ff5a3d2771008385f636d1f79 ===
% 用户上传的 Irodov Arihant 底本，PDF 页88--89，原书印刷页79--80。
% 原题 146--159；原书图2.4两种T-S循环已经对照图像检查。
\thermoentropyproblem{146}{79}
一摩尔定容摩尔热容为 $C_V$ 的理想气体经历一个过程，
在过程中其熵与温度满足 $S=\alpha/T$，
其中 $\alpha$ 为常数。气体温度由 $T_1$ 变化至 $T_2$。
求：(a) 过程摩尔热容 $C(T)$；
(b) 传入气体的热量 $Q$；
(c) 气体做功 $A$。
\exsolution
过程摩尔热容的定义为
$C(T)=T\,\mathrm dS/\mathrm dT$。
从 $S=\alpha/T$ 求导得
\[
\boxed{C(T)=-\frac{\alpha}{T}}.
\]
沿准静态过程
$\delta Q=T\,\mathrm dS=-\alpha\,\mathrm dT/T$，
故
\[
\boxed{Q=-\alpha\int_{T_1}^{T_2}\frac{\mathrm dT}{T}
=\alpha\ln\frac{T_1}{T_2}}.
\]
一摩尔理想气体的内能变化
$\Delta U=C_V(T_2-T_1)$，
热力学第一定律 $Q=\Delta U+A$ 给出
\[
\boxed{A=\alpha\ln\frac{T_1}{T_2}
+C_V(T_1-T_2)}.
\]
\excheck $C(T)$ 是\emph{过程}热容而非 $C_V$；
只要 $T>0$，其符号取决于 $\alpha$，
热量和外功都必须保留代数符号。

\thermoentropyproblem{147}{79}
某工质在图2.4(a)、(b) 所示的两种
$T$--$S$ 三角形循环中工作，
循环最高温与最低温之比为
$n=T_h/T_c>1$。求两种循环的热效率。
\exsolution
在温度—熵图上，沿可逆支路的热量
$\delta Q=T\,\mathrm dS$；
竖直支路熵恒定，不交换热。
设两端水平熵间隔为 $\Delta S>0$。
(a) 按原图箭头，顶部恒温支路向右，
吸热
$Q_{\rm in}=T_h\Delta S$；
倾斜支路返回时向外放热，
由于直线在 $T$--$S$ 图上的面积为梯形，
放热模为
$Q_{\rm out}=(T_h+T_c)\Delta S/2$。
所以
\[
\boxed{\eta_a=1-\frac{Q_{\rm out}}{Q_{\rm in}}
=1-\frac{T_h+T_c}{2T_h}
=\frac{n-1}{2n}}.
\]
(b) 原图沿上升斜线向右吸热，
$Q_{\rm in}=(T_h+T_c)\Delta S/2$；
底部等温线向左排热
$Q_{\rm out}=T_c\Delta S$，
因此
\[
\boxed{\eta_b=1-\frac{2T_c}{T_h+T_c}
=\frac{n-1}{n+1}}.
\]
\excheck 两循环都是顺时针正功方向；
必须依据原图箭头判别吸热段，
而不能只看两个三角形的外形。
两个效率均小于卡诺上界 $1-1/n$。

\thermoentropyproblem{148}{79}
两个彼此绝热的容器用带阀门的短管相连，
其中一个装有 $\nu=2.2$ mol 理想气体，
另一个抽成真空。阀门打开后，
气体的可及体积增大为原来的 $n=3.0$ 倍。
求气体的熵增。
\exsolution
将两容器整体视为隔热、刚性系统。
气体向真空自由膨胀，没有外功也没有热交换，
因此 $\Delta U=0$。
理想气体内能只依赖温度，
故末态温度等于初温。
熵为状态函数，可以用假想可逆等温膨胀
连接相同初末态：
\[
\Delta S=\int_{V_1}^{V_2}\frac{\nu R}{V}\,\mathrm dV
=\nu R\ln(V_2/V_1).
\]
于是
\[
\boxed{\Delta S=\nu R\ln n
=(2.2)(8.314)\ln3
\simeq20.1\,\mathrm{J/K}}.
\]
\excheck 实际自由膨胀不可逆，不能沿实际路径写
$\delta Q_{\rm actual}=T\,\mathrm dS$；
熵积分使用的是假想可逆路径。

\thermoentropyproblem{149}{79}
一个绝热圆筒被无质量活塞分成体积相等的两半。
左半部装有一摩尔绝热指数为 $\gamma$、
初温 $T_0$ 的理想气体，右半部是真空。
释放活塞，使气体充满整个圆筒。
随后缓慢将活塞移回原位置。
求这两个过程完成之后气体的内能增量和熵增。
\exsolution
释放时气体向真空自由膨胀到原体积的两倍；
由于整体绝热、无外功，
$\Delta U_1=0$ 且 $T_1=T_0$。
自由膨胀产生熵
$\Delta S_1=R\ln2$。
第二步缓慢移动活塞，
把气体由 $2V_0$ 压回 $V_0$。
按题设圆筒隔热与缓慢准静态运动，
这是可逆绝热压缩；
由 $TV^{\gamma-1}=\text{常数}$，
\[
T_f=T_0\left(\frac{2V_0}{V_0}\right)^{\gamma-1}
=2^{\gamma-1}T_0.
\]
一摩尔总内能增加
\[
\boxed{\Delta U=C_V(T_f-T_0)
=\frac{RT_0}{\gamma-1}(2^{\gamma-1}-1)}.
\]
第二步熵不变，故总熵增为
\[
\boxed{\Delta S=R\ln2}.
\]
\excheck 原书答案文本抽取可能把
$2^{\gamma-1}$ 显示成“$2\gamma-1$”；
由绝热温度体积关系确认指数结构。

\thermoentropyproblem{150}{79}
同一初始态的理想气体在不与外界交换热量的条件下
膨胀到相同的终体积 $V$。
比较：(a) 快速膨胀；(b) 极其缓慢膨胀。
末态压强是否相同？
\exsolution
对同一初态、相同终体积，
如果慢膨胀始终可逆绝热，
末态满足 $pV^\gamma=\text{常数}$。
而实际快速膨胀通常产生不可逆耗散，
同一终体积下其熵高于可逆绝热终态。
理想气体的熵差可写为
\[
\Delta S=C_V\nu\ln(T_f/T_i)
+\nu R\ln(V_f/V_i).
\]
两过程的初态与终体积相同，
较高的终态熵对应较高的终态温度。
由 $p_f=\nu RT_f/V_f$，
得出
\[
\boxed{p_{f,\rm fast}>p_{f,\rm slow}}
\]
（快速过程出现熵产生、而慢过程可逆时）。
\excheck 这是一项\emph{关于实际快过程与可逆慢过程}
的比较：如果两个过程具有相同的实际功和能量耗散，
就不能仅凭“快慢”二字确定严格的压强差。

\thermoentropyproblem{151}{80}
一个绝热容器被隔板分成两部分，
大的一侧体积是小的一侧的 $n=2.0$ 倍。
小室内有 $\nu_1=0.30$ mol 氮气，
大室内有 $\nu_2=0.70$ mol 氧气。
两气体初温相同。
在隔板上穿孔让两气体混合，
视为理想气体，求系统熵增。
\exsolution
设小室容积为 $V$、大室为 $nV$。
无热交换、无体积功且两气体初温相同，
理想气体内能守恒表明末温不变。
每种气体可等温地从各自初始容积
扩散到共同总容积 $(n+1)V$。
因此
\[
\Delta S_1=\nu_1R\ln\frac{(n+1)V}{V},\qquad
\Delta S_2=\nu_2R\ln\frac{(n+1)V}{nV}.
\]
相加：
\[
\boxed{\Delta S=
\nu_1R\ln(1+n)
+\nu_2R\ln(1+1/n)}.
\]
取给定数值，
\[
\boxed{\Delta S=R(0.30\ln3+0.70\ln1.5)
\simeq5.10\,\mathrm{J/K}}.
\]
\excheck 氮、氧是不同分子，
均可占据最终整个容积；
不能将混合熵误写成同一种气体移除隔板的零熵变。

\thermoentropyproblem{152}{80}
把质量 $m_1=300\,\mathrm g$、
温度 $t_1=97^\circ\mathrm C$ 的铜块，
投入质量 $m_2=100\,\mathrm g$、
温度 $t_2=7^\circ\mathrm C$ 的水中。
假定热量计热容可忽略，整体隔热。
求达到温度平衡时系统的总熵增。
\exsolution
取铜和水的质量比热分别为
$c_1,c_2$，在本温区视为常数。
能量守恒：
\[
m_1c_1(T_f-T_1)+m_2c_2(T_f-T_2)=0,
\quad
T_f=\frac{m_1c_1T_1+m_2c_2T_2}
{m_1c_1+m_2c_2}.
\]
各自熵变由可逆积分求得：
\[
\boxed{\Delta S=
m_1c_1\ln\frac{T_f}{T_1}
+m_2c_2\ln\frac{T_f}{T_2}}.
\]
用 $c_1\simeq0.39\,\mathrm{J/(g\,K)}$、
$c_2\simeq4.18\,\mathrm{J/(g\,K)}$，
$T_1=370.15\,\mathrm K$、
$T_2=280.15\,\mathrm K$，
先算得 $T_f\simeq300\,\mathrm K$，
再代熵式：
\[
\boxed{\Delta S\simeq3.7\,\mathrm{J/K}}
\quad(c_{\rm Cu}=0.39,\ c_w=4.18\ \mathrm{J/(g\,K)}).
\]
\excheck 熵变应使用\emph{开尔文温度}取对数，
不能将摄氏温度比值直接代入。
本版书末给约 $4.4\,\mathrm{J/K}$；
该数值依赖所用铜、水比热表和温度积分近似，
本稿明确列出所用常数与公式，不直接照搬近似数值。

\thermoentropyproblem{153}{80}
两只相同体积、彼此绝热的刚性容器，
各装一摩尔\emph{同一种}理想气体，
初温分别为 $T_1,T_2$。
其定容摩尔热容为 $C_V$。
打开阀门后两容器气体达到共同平衡，
求气体总熵增并证明非负。
\exsolution
整体绝热且体积不变，
相同气体每容器各一摩尔，
故能量守恒给
$C_V(T_1+T_2)=2C_VT_f$，
\[
T_f=\frac{T_1+T_2}{2}.
\]
按各自一摩尔的熵变加和；
两容器总体积相同、最终混合的气体为同一种，
若用分室可逆路径计算，体积项总体抵消，
因而
\[
\Delta S=C_V\ln(T_f/T_1)
+C_V\ln(T_f/T_2)
=\boxed{C_V\ln\frac{(T_1+T_2)^2}{4T_1T_2}}.
\]
由算术—几何平均不等式
$(T_1+T_2)^2\ge4T_1T_2$，
所以
\[
\boxed{\Delta S\ge0},
\]
且仅在 $T_1=T_2$ 时等号成立。
\excheck 两种气体若是不同种类，还需包括混合熵；
本题明确为\emph{同一种气体}，不可加虚假混合熵。

\thermoentropyproblem{154}{80}
有 $N$ 个氦原子处于室温、体积为
$1.0\,\mathrm{cm^3}$ 的立方容器中。
求：(a) 所有原子在同一瞬间全部聚集于一半容积的概率；
(b) 若要求这一事件在约
$t=10^{10}$ 年内发生一次，需要大致多少个原子？
\exsolution
若原子位置独立且在两半室等概率分布，
指定某一半的每个原子都位于其中的概率为 $1/2$，
故
\[
\boxed{P=2^{-N}}.
\]
(b) 设重新抽样出近似独立位置分布的时间尺度为
$\tau$，若观察时长为 $t$，
估计独立尝试次数为 $K\simeq t/\tau$。
每次概率为 $2^{-N}$，
平均发生次数 $K2^{-N}\sim1$，故
\[
\boxed{N\simeq\frac{\ln(t/\tau)}{\ln2}}.
\]
取室温原子通过厘米量级容器的时间
$\tau\simeq10^{-5}\,\mathrm s$、
$t\simeq3.16\times10^{17}\,\mathrm s$，
得到
$\boxed{N\simeq75\ \text{个，约80个}}$。
\excheck (b) 只是基于近似独立尝试的数量级估计；
真实微观位置关联时间不严格等于一次过箱时间。
(a) 按原书是\emph{指定}的一半区域，
若接受左右任一半，概率须改为 $2^{1-N}$。

\thermoentropyproblem{155}{80}
把 $N=10$ 个相同分子分配到圆筒等体积的两半，
求最概然分配的统计权重及其发生概率。
\exsolution
各分子可独立位于左半或右半，
宏观态“左半有 $k$ 个”对应的微观排列数为
\[
\Omega_k=\binom Nk=\frac{N!}{k!(N-k)!}.
\]
对于偶数 $N=10$，最概然数为 $k=N/2=5$，
所以
\[
\boxed{\Omega_{\rm max}=\binom{10}{5}
=\frac{10!}{(5!)^2}=252}.
\]
共有 $2^{10}$ 个等概率左右分配，
概率为
\[
\boxed{P_{\rm max}=\frac{252}{2^{10}}
=\frac{252}{1024}\simeq24.6\%}.
\]
\excheck “分子相同”指它们物理种类相同；
排列计数仍可以对可追踪的微观粒子位置计数，
不应把所有分配只算成一个状态。

\thermoentropyproblem{156}{80}
容器装有 $N$ 个理想气体分子，
把容积设想划为相等的 $A,B$ 两半。
求 $A$ 半中有恰好 $n$ 个分子的概率。
特别列出 $N=5$ 时
$n=0,1,2,3,4,5$ 的概率。
\exsolution
每个分子进入 $A$ 的概率为 $1/2$，
所有分子位置独立时，
恰好 $n$ 个在 $A$ 的概率服从二项分布：
\[
\boxed{P_n=\binom Nn
\left(\frac12\right)^n
\left(\frac12\right)^{N-n}
=\frac{N!}{n!(N-n)!\,2^N}}.
\]
$N=5$ 的六个概率依次为
\[
\boxed{
P_0=\frac1{32},\
P_1=\frac5{32},\
P_2=\frac{10}{32},\
P_3=\frac{10}{32},\
P_4=\frac5{32},\
P_5=\frac1{32}}.
\]
\excheck 六项相加
$(1+5+10+10+5+1)/32=1$，
满足概率归一化。

\thermoentropyproblem{157}{80}
体积为 $V_0$ 的容器内有 $N$ 个理想气体分子。
任选其中体积为 $V$ 的固定子区域，
求其中恰有 $n$ 个分子的概率，
并讨论 $V=V_0/2$。
\exsolution
位置均匀分布时，单个分子落入子区域的概率
$q=V/V_0$。
恰有 $n$ 个分子落入的组合数为 $\binom Nn$，
因此
\[
\boxed{P_n=\binom Nn
\left(\frac V{V_0}\right)^n
\left(1-\frac V{V_0}\right)^{N-n}}.
\]
若 $V=V_0/2$，
则 $q=1/2$，
\[
\boxed{P_n=\binom Nn2^{-N}}.
\]
\excheck 需要 $0\le V\le V_0$ 与 $0\le n\le N$；
所有 $n$ 的概率之和由二项式定理等于1。

\thermoentropyproblem{158}{80}
标准状况下的理想气体，
求一个球形取样区域的直径 $d$，
使区域内分子数的相对涨落为
$\eta=1.0\times10^{-3}$。
同时求该区域内平均分子数。
\exsolution
对均匀稀薄气体中的固定小取样区域，
当整体系统远大于该区域，
分子数涨落近似泊松分布：
$\operatorname{Var}N=\bar N$。
相对标准差为
\[
\eta=\frac{\sqrt{\operatorname{Var}N}}{\bar N}
=\frac1{\sqrt{\bar N}},
\]
故
\[
\boxed{\bar N=\eta^{-2}=10^6}.
\]
标准状况数密度
$n_0=p_0/(k_BT_0)
\simeq2.69\times10^{25}\,\mathrm{m^{-3}}$。
球体积
$\frac{\pi d^3}{6}=\bar N/n_0$，
所以
\[
\boxed{d=\left(\frac6{\pi n_0\eta^2}\right)^{1/3}
\simeq4.1\times10^{-7}\,\mathrm m
=0.41\,\mathrm{\mu m}}.
\]
\excheck 由体积求直径必须取\emph{三次方根}；
PDF 提取文本有时丢掉此根号结构。

\thermoentropyproblem{159}{80}
一摩尔单原子理想气体装在等容容器中，
初温 $T_0=300\,\mathrm K$，
加热 $\Delta T=1.0\,\mathrm K$。
求气体统计权重 $\Omega$ 变化的倍率。
\exsolution
玻尔兹曼熵公式
$S=k_B\ln\Omega$，
故统计权重比
$\Omega_f/\Omega_i=e^{\Delta S/k_B}$。
等容加热一摩尔单原子气体的熵增：
\[
\Delta S=C_V\ln\frac{T_0+\Delta T}{T_0},
\quad C_V=\frac32R=\frac32N_Ak_B.
\]
因此
\[
\boxed{\frac{\Omega_f}{\Omega_i}
=\left(1+\frac{\Delta T}{T_0}\right)^{3N_A/2}}.
\]
用十进制数量级表示：
\[
\log_{10}\frac{\Omega_f}{\Omega_i}
=\frac{3N_A}{2}\log_{10}\left(\frac{301}{300}\right)
\simeq1.31\times10^{21},
\]
即
\[
\boxed{\Omega_f/\Omega_i
\simeq10^{\,1.31\times10^{21}}}.
\]
\excheck 这是极巨大的\emph{无量纲倍率}，
不能误写为简单的 $10^{21}$。
% === END 12-entropy-cycles-146-159.tex ===
